%%%%%%%%%%%%%%%%%%%%%%%%%%%%%%%%%%% Einstellungen/Vorgaben %%%%%%%%%%%%%%%%%%%%%%%%%%%%%%%%%%%

% Basics
\documentclass[ngerman, a4paper, 12pt, oneside, pdftex]{report}

% UTF-8 codierung verwenden
\usepackage[utf8]{inputenc}

\clubpenalty = 10000
\widowpenalty = 10000
\displaywidowpenalty = 10000

% Nützliche Packages
\usepackage[T1]{fontenc}
\usepackage{lmodern} % Vektor-Schriftart
\usepackage{geometry} % Seitenränder
\usepackage{graphicx} % Grafiken
\usepackage{caption} % Bildunterschriften etc.
\usepackage{acronym} % Abkürzungen
\usepackage{listings} % Listenverzeichnis
\usepackage{color} % Benutzerdefinierte Farben
\usepackage{csquotes}
\usepackage{svg} % Vektor-Grafiken
\usepackage{enumitem}
\usepackage{tcolorbox}
\usepackage[htt]{hyphenat}
\usepackage{scrextend}
\usepackage[bottom]{footmisc}
\usepackage[german]{babel}
%\usepackage[backend=biber,maxcitenames=3,hyperref=true,bibencoding=inputenc,style=alphabetic,citestyle=alphabetic]{biblatex} % Literaturverzeichnis
\usepackage{setspace}
%\usepackage[nottoc, notlof, notlot]{tocbibind}
\usepackage[headsepline]{scrpage2} % Kopf- und Fusszeilen

% Farben
\definecolor{dkgreen}{rgb}{0,0.6,0}
\definecolor{gray}{rgb}{0.5,0.5,0.5}
\definecolor{mauve}{rgb}{0.58,0,0.82}

% Code-Embedding
\lstset{frame=single,
  language=Java,
  aboveskip=3mm,
  belowskip=3mm,
  showstringspaces=false,
  columns=flexible,
  basicstyle={\small\ttfamily},
  numbers=left,
  numberstyle=\tiny\color{gray},
  keywordstyle=\color{blue},
  commentstyle=\color{dkgreen},
  stringstyle=\color{mauve},
  keepspaces=true,
  breaklines=true,
  breakatwhitespace=true,
  tabsize=3
}

% Sprache
\selectlanguage{german}

% Seitenränder
\geometry{left=25mm, right=25mm, top=25mm, bottom=25mm} % bindingoffset=5mm

\setlength{\headheight}{1.1\baselineskip}

% Zeilenabstand 1.5
\onehalfspacing

% Caption Setup
\captionsetup{labelfont=bf, textfont=normalfont}

% Verfasser
\newcommand{\Autor}{Jakob Gietl und Fynn Arnold}
\newcommand{\Kurs}{TINF17B3}

% Arbeit
\newcommand{\ArbeitTitel}{Implementierung eines PIC-Simulators in Java}
\newcommand{\ArbeitZeitraum}{?? Wochen}
\newcommand{\ArbeitAbgabeDatum}{??.??.????}

% Hochschule
\newcommand{\DHName}{Duale Hochschule Baden-Württemberg}
\newcommand{\DHStadt}{Karlsruhe}
\newcommand{\DHStudiengang}{Informationstechnik}
\newcommand{\DHTitelseiteLogo}{\includegraphics[width=4cm]{img/dhbw-logo.png}}

% Eigene Befehle
\newcommand{\todo}[1]{\begin{tcolorbox}[colback=green,colframe=black,boxrule=1pt,left=0pt,right=0pt,top=0pt,bottom=0pt,sharp corners]\small{\textit{\textbf{TODO:} #1}}\end{tcolorbox}}
\newcommand{\descr}[2]{\noindent \textbf{#1} \begin{addmargin}[10mm]{0mm} #2 \end{addmargin} \vskip 12pt}
\newcommand{\source}[1]{\caption*{\textbf{Quelle:} #1}}
   
% Line Breaks
%\def\UrlBreaks{\do\/\do-}
%\setcounter{biburllcpenalty}{9000}% Kleinbuchstaben
%\setcounter{biburlucpenalty}{9000}% Großbuchstaben

% Literaturverzeichnis
%\addbibresource{./Anhang/Literaturverzeichnis.bib}

\usepackage[pdftex,
            pdfauthor={\Autor},
            pdftitle={\ArbeitTitel}]{hyperref}

% Kopfzeile
\pagestyle{scrheadings}
\clearscrheadfoot

\setheadsepline{0.4pt}
\automark[section]{chapter}
\ihead{\headmark}
\ohead{\leftmark}

\newcounter{savepage} % Seitenzahl

\begin{document}



%%%%%%%%%%%%%%%%%%%%%%%%%%%%%%%%%%% Titelseite %%%%%%%%%%%%%%%%%%%%%%%%%%%%%%%%%%%

\hypersetup{pageanchor=false}

\begin{onehalfspacing}
\begin{titlepage}
\begin{center}
\vspace*{-2cm}
\hfill
\DHTitelseiteLogo \\
\hfill
\\[40mm]
{\Huge \mbox{\ArbeitTitel}}\\[20mm]
{\large \DHName \ \DHStadt}\\[5mm]
{\large Studiengang \DHStudiengang}\\[20mm]
{\large\bfseries \Autor}\\[5mm]
{\large Kurs \Kurs}
\vfill
\end{center}

\begin{tabular}{l@{\hspace{15mm}}l}
Abgabedatum		& \ArbeitAbgabeDatum \\
Bearbeitungszeitraum	& \ArbeitZeitraum
\end{tabular}
\end{titlepage}
\end{onehalfspacing}



%%%%%%%%%%%%%%%%%%%%%%%%%%%%%%%%%%% Verzeichnisse %%%%%%%%%%%%%%%%%%%%%%%%%%%%%%%%%%%

\clearpage
\cfoot[\pagemark]{\pagemark} % Fußzeile mit Seitennummer
\pagenumbering{Roman} % Römische Nummerierung

\begin{singlespace}

% Inhaltsverzeichnis
\tableofcontents
\noindent Hier wird das Inhaltsverzeichnis stehen!

% Abbildungsverzeichnis
\listoffigures
\addcontentsline{toc}{chapter}{Abbildungsverzeichnis}
\noindent Hier wird das Abbildungsverzeichnis stehen!

% Tabellenverzeichnis
\listoftables
\addcontentsline{toc}{chapter}{Tabellenverzeichnis}
\noindent Hier wird das Tabellenverzeichnis stehen!


\end{singlespace}

% Quellcodeverzeichnis
\renewcommand{\lstlistlistingname}{Quellcodeverzeichnis}
\lstlistoflistings % Quellcodeverzeichnis
\addcontentsline{toc}{chapter}{Quellcodeverzeichnis}


\newpage
\hypersetup{pageanchor=true}



%%%%%%%%%%%%%%%%%%%%%%%%%%%%%%%%%%% Kapitel %%%%%%%%%%%%%%%%%%%%%%%%%%%%%%%%%%%

\setcounter{savepage}{\value{page}}
\pagenumbering{arabic}

\LaTeX

% Kapitel 1: Einleitung
%\input{chapter1.tex}

% Kapitel 2: Grundlagen
%\input{chapter2.tex}

% Kapitel 3: Planung
%\input{chapter3.tex}

% Kapitel 4: Realisierung
%\input{chapter4.tex}

% Kapitel 5: Fazit
%\input{chapter5.tex}



%%%%%%%%%%%%%%%%%%%%%%%%%%%%%%%%%%% Literaturverzeichnis %%%%%%%%%%%%%%%%%%%%%%%%%%%%%%%%%%%

%\newpage
%\pagestyle{empty}
%\pagenumbering{Roman} \setcounter{page}{\value{savepage}}


%\renewcommand{\bibname}{Literaturverzeichnis}
%\addcontentsline{toc}{chapter}{\bibname}
%\printbibliography
%\noindent Hier wird das Literaturverzeichnis (auch mit Titel) stehen!

%\newpage\null\thispagestyle{empty}\newpage % leere Seite



%%%%%%%%%%%%%%%%%%%%%%%%%%%%%%%%%%% Ende des Dokuments %%%%%%%%%%%%%%%%%%%%%%%%%%%%%%%%%%%

\end{document}
